% Teaching example: F=diag(100,1), g=(1,1), epsilon=0.005.
% Both axes use 240 mm per parameter unit; the ellipse semiaxes are 0.01,0.1.
% The equal-budget endpoints are computed from the formulas in the text.
% Keep the narrow ellipse and annotate from its sides; blue/dashed/circle denotes
% Euclidean descent, red/solid/square denotes natural descent.
\begin{tikzpicture}[
  x=1mm,y=1mm,
  every node/.style={font=\figurefont\footnotesize,text=FigureInk,inner sep=1pt},
  >={Latex[length=1.5mm,width=1.1mm]}
]
  \path[use as bounding box] (0,0) rectangle (112,66);
  \coordinate (O) at (52,34);
  \coordinate (E) at ({52-24/sqrt(101)},{34-24/sqrt(101)});
  \coordinate (F) at ({52-0.24/sqrt(1.01)},{34-24/sqrt(1.01)});

  \fill[FigurePaper] (O) ellipse [x radius=2.4mm,y radius=24mm];
  \draw[FigureMuted,line width=.8pt] (O) ellipse [x radius=2.4mm,y radius=24mm];
  \draw[->,FigureStroke,line width=.6pt] (26,34)--(77,34) node[right] {\(\delta_1\)};
  \draw[->,FigureStroke,line width=.6pt] (52,6)--(52,62) node[above] {\(\delta_2\)};
  \draw[FigureStroke,line width=.6pt] (51,58)--(53,58);
  \draw[FigureStroke,line width=.6pt] (51,10)--(53,10);
  \node[anchor=west] at (56,58) {\(0.1\)};
  \node[anchor=west] at (56,7) {\(-0.1\)};
  \node[anchor=south west] at (54,35) {\(0\)};
  \node[anchor=east] at (45,49) {\(\delta_1=\pm0.01\)};
  \draw[FigureMuted,line width=.5pt] (45,46)--(49.6,34);

  \draw[->,FigureBlue,dashed,line width=1pt] (O)--(E);
  \fill[FigureBlue] (E) circle (.7mm);
  \draw[FigureBlue,line width=.6pt] (E)--(36,26)--(22,26);
  \node[anchor=east,text=FigureBlue] at (34,22) {欧氏方向\(\bm\delta_{\mathrm E}\)};

  \draw[->,FigureRed,line width=1pt] (O)--(F);
  \fill[FigureRed] ($(F)+(-.7,-.7)$) rectangle ($(F)+(.7,.7)$);
  \draw[FigureRed,line width=.6pt] (F)--(66,16);
  \node[anchor=west,text=FigureRed] at (68,16) {自然方向\(\bm\delta_{\mathrm F}\)};
\end{tikzpicture}
